The reconstruction of a smooth field onto a fixed grid is a necessary step for direct comparisons to various real-world observations. Projecting SPH data onto a fixed grid becomes challenging in adaptive environments, where some particles may have smoothing lengths far below the grid size, whilst others are resolved by thousands of pixels. In this paper we show how the common approach of treating particles below the grid size as Monte Carlo tracers of the field leads to significant reconstruction errors, and despite good convergence properties is unacceptable for use in synthetic observations in astrophysics. We propose a new method, where particles smaller than the grid size are `blitted' onto the grid using a high-resolution pre-calculated kernel, and those close to the grid size are subsampled, that allows for converged predictions for projected quantities at all grid sizes.